%% 
%% Copyright 2019-2024 Elsevier Ltd
%% 
%% Version 2.4
%% 
%% This file is part of the 'CAS Bundle'.
%% --------------------------------------
%% 
%% It may be distributed under the conditions of the LaTeX Project Public
%% License, either version 1.2 of this license or (at your option) any
%% later version.  The latest version of this license is in
%%    http://www.latex-project.org/lppl.txt
%% and version 1.2 or later is part of all distributions of LaTeX
%% version 1999/12/01 or later.
%% 
%% The list of all files belonging to the 'CAS Bundle' is
%% given in the file `manifest.txt'.
%% 
%% Template article for cas-dc documentclass for 
%% double column output.

%\documentclass[a4paper,fleqn,longmktitle]{cas-dc}
\documentclass[a4paper,fleqn]{cas-dc}

%\usepackage[authoryear,longnamesfirst]{natbib}
%\usepackage[authoryear]{natbib}
\usepackage[numbers]{natbib}

\usepackage{algorithm}
\usepackage[noend]{algpseudocode}

%%%Author definitions
\def\tsc#1{\csdef{#1}{\textsc{\lowercase{#1}}\xspace}}
\tsc{RNG}
%\tsc{WGM}
%\tsc{QE}
%\tsc{EP}
%\tsc{PMS}
%\tsc{BEC}
%\tsc{DE}
%%%

\begin{document}

\let\WriteBookmarks\relax
\def\floatpagepagefraction{1}
\def\textpagefraction{.001}
\shorttitle{Expected Linear Time Algorithm for Computing 3D Relative Neighborhood Graphs}
\shortauthors{Z. Zzyzek}

\title [mode = title]{Expected Linear Time Algorithm for Computing 3D Relative Neighborhood Graphs}                      

%\tnotemark[1,2]
%\tnotetext[1]{This document is the results of the research project funded by the National Science Foundation.}
%\tnotetext[2]{The second title footnote which is a longer text matter to fill through the whole text width and overflow into another line in the footnotes area of the first page.}


\author[1]{Z. Zzyzek}[role=Researcher,
                   orcid=0009-0005-2175-1166]
\cormark[1]
\fnmark[1]
\ead{zzyzek@thumpernet.com}
\ead[url]{zzyzek.github.io}

%\credit{Conceptualization of this study, Methodology, Software}

\begin{abstract}
  An expected linear time algorithm for computing the relative neighborhood graph (RNG) for points in Euclidean space,
  distributed uniformly in a unit three dimensional cube, is given.
  For each point, potential neighbors for an RNG connection are limited to an expected small bounding volume.
  Volumes relative to the considered point are bounded through half plane partitions,
  constructed from the point and its neighbors,
  that remove regions where no RNG edge connections to the point can occur.
  %The bounding volume is found by fencing in enclosing a convex hull 
  %where the bounding volume is found through a plane partition heuristic.
  The small bounding volume provides an expected constant number of neighbors to consider, making
  each single point RNG calculation $O(1)$, providing the $O(n)$ expected runtime when enumerating through all points.
  %The $K$ surrounding potential neighbors are then tested with
  %a naive $O(K^3)$ algorithm.
  %The RNG has finite degree for random point distributions in three dimensions,
  %resulting in an expected finite $K$, making $O(K^3) = O(1)$, providing 
  %the linear expected time algorithm when enumerating through all points.
  To the author's knowledge, this improves on the previously best known $O(n \log n)$ in three dimensions.
\end{abstract}

%%% CREATES INITIAL 'GRAPHICAL ABSTRACT' PAGE
%%%
%\begin{graphicalabstract}
%\includegraphics{figs/cas-grabs.pdf}
%\end{graphicalabstract}

%%% CREATES INITIAL 'HIGHLIGHTS' PAGE
%%%
%\begin{highlights}
%\item Research highlights item 1
%\item Research highlights item 2
%\item Research highlights item 3
%\end{highlights}

\begin{keywords}
  graph algorithm \sep relative neighborhood graph \sep computation geometry \sep RNG
%quadrupole exciton \sep polariton \sep \WGM \sep \BEC

\end{keywords}


\maketitle

\section{Introduction}

Let $P = (p _ 0, p _ 1, \dots, p _ {n-1})$ be a set of points in $\mathbb{R}^3$.
The relative neighborhood graph, $\text{RNG}(P)$, is an undirected, simple graph that connects two vertices,
$p _ i, p _ j$ if and only if there is no other vertex within the \emph{lune} between
them.

$$
\begin{array}{l}
  (p _ i, p _ j) \in \text{RNG}(P) \\
  \iff |p _ i - p _ j| \le \max _ { k \ne i,j } ( |p _ i - p _ k|, |p _ j - p _ k| )
\end{array}
$$

Here, $|\cdot|$ is the standard Euclidean norm.
This paper focuses almost exclusively on points in 3D but will occasionally talk about
points in 2D for explanatory examples.

Figure \ref{TK} shows a highlighted 2D example of a connections and precluded connection
and Figure \ref{TK}, \ref{TK} shows examples of a relative neighborhood graph for 2D and 3D.

When points are placed uniformly at random in the unit cube, the maximum
degree of the relative neighborhood graph is bounded by a constant.
When random in this manner, the edges from any point connecting it to its relative
neighborhood graph only extend to a proportionally small local neighborhood.
This local neighborhood is effectively constant in neighbor size, allowing us to calculate
each points relative neighborhood graph edges in constant time, providing a linear time algorithm
to find the relative neighborhood graph in 3D.


It's believed that the worst case runtime for a relative neighborhood graph in 2D and 3D is
at least $O(n \log n)$.
Randomness holds special structure and many cases in the general case won't happen, with
high probability.
Note that randomness is a different condition than generic or general position, as point
configurations that pass a generic or general position criteria might not be considered
random.
Only points uniformly distributed in the unit cube are considered in this paper,
without focusing on whether they're in generic or general position.

%Call the relative neighborhood graph around a point to be $\text{RNG}(P,p)$.
%That is, the subset of the relative neighborhood graph that only has an edge that includes $p$.





\section{Related Work}

OUTLINE

- Talk about KN85 and KNT86 and their exceptions to corner and side regions of grid.
  Also talk about restriction to 2d and how generalizations to 3d are complicated

- 3d algs of agarwal,matousek are meant for worst case setups but super linear (e.g. advesarial)

- talk about efficient computation of Urquhart (sp?) graph

- talk about using delaunay triangulation to bootstrap

- review gives best known algorithm at $O(n \log n)$ for general position (doesn't talk about random position?)

- Jaromczyk and Kowaluk on $O(n^2 \log n)$ 3d algorithm with 'coarse' elimination phase that was inspiration
  for plane cut heuristic


\section{Algorithm}

All points eventually will be evaluated through a naive relative neighborhood graph
calculation, relative to an \emph{anchor point}, $p$.
Algorithm \ref{alg:RNGPointNaive} provides the pseudo code that does the all pairs
test of points in $Q$ to the anchor point, $p$.

Algorithm \ref{alg:RNGPointNaive} has $O(|Q|^2)$ runtime.
As will be detailed
in a later section, $|Q|$ will be shown to be almost surely bounded above by a constant, making the naive
pairwise point testing of Algorithm \ref{alg:RNGPointNaive} constant time.

\begin{algorithm}
  \caption{Single point naive Relative Neighborhood Graph for the list $Q$ to test for connections to $p$}\label{alg:RNGPointNaive}
  \begin{algorithmic}[1]
    \Procedure{RNGPointNaive}{Q, p}
      \For{ $q \in Q / \{p\}$ }
        \State $\text{isConnected } = \text{ \emph{False}}$
        \For{ $u \in Q / \{p,q\}$ }
          \If{ $\text{InLune}(p,q,u)$ }
            \State $\text{isConnected } = \text{ \emph{False}}$
          \EndIf
        \EndFor
        \If{ $\text{isConnected } == \text{ \emph{True}}$ }
          $E = E \cup (p,q)$
        \EndIf
      \EndFor

      \State $\text{return } E$
    \EndProcedure
  \end{algorithmic}
\end{algorithm}

Algorithm \ref{alg:RNGSPoIF} intializes the grid and then processes
each point individually with a call to the Algorithm \ref{alg:RNGpSPoIF}.

\begin{algorithm}
  \caption{Relative Neighborhood Graph, Secure Posts on Increasing Fence (SPoIF)}\label{alg:RNGSPoIF}
  \begin{algorithmic}[1]
    \Procedure{RNGSPoIF}{P}
    \State \vspace{0.5em} $E = \{\}$
    \State \vspace{0.5em} $N = |P|, \ \ \ \ N_s = \lceil N^{\frac{1}{3}} \rceil, \ \ \ \ s = \frac{1}{N_s}$,
    \State \vspace{0.5em} $B[0:N_s][0:N_s][0:N_s] = \emptyset$
    \For{$p \in P$}
      \vspace{0.5em} \State $\text{append}( B[ \lfloor N_s  p_x \rfloor ][ \lfloor N_s  p_y \rfloor ][ \lfloor N_s  p_z \rfloor ] , p )$
    \EndFor

    \For{$p \in P$}
      \vspace{0.5em} \State $E = E \cup \text{RNGPointSPoIF}(B,P,p)$
    \EndFor

    \State $\text{return } E$

    \EndProcedure
  \end{algorithmic}
\end{algorithm}

Algorithm \ref{alg:RNGpSPoIF} processes each point individually by
looking at a local grid neighborhood of the grid cell occupied by the current
anchor point $p$ being processed.
Points within the neighboring cells of $p$
are then tested to see if they've secured the fence.

Conceptually, cutting planes are built that pass through neighbor points, $q$, with
normal in the direction of $p$ to $q$.
The cutting planes are tested against a bounding cube to see if they partition
the edge faces from $p$.

The bounding cube faces are called a \emph{fence}.
A fence face that has been partitioned from the anchor point $p$ is said to be \emph{secured},
with the whole fence being secured once all fence faces have been secured.

To test to see if the fence is secured, points on the fence face are used
as representatives, which will be called \emph{fence posts}.
Small patches of the fence face are then represented by clusters of fence posts,
where the fence posts create the bounding vertices of the square.
The clusters of fence posts can be tested so that, if they are partitioned
from $p$, this will ensure that the patch the fence posts bound is also partitioned.

Fence posts clusters that have been partitioned, or secured, imply
fence patches have been secured.
Once all patches have been secured, the whole fence will be secured.
The secured fence will then be used to collect nearby points to be
passed to a naive relative neighborhood graph calculation relative
to the anchor point, $p$.

% WIP!!!
A fence face patch square, reprsented by a cluster of fence posts at the corners,
is tested against a plane that passes through a point $q$ and has normal
in the anchor point $p$
of the fence face pathc

Fence securitization is done by testing all fence post points in the frusturm that
make the fence face patch, called \emph{fence posts}, are \emph{secured}.
The vector in the direction of $(p-q)$ and the fence post point relative
to the grid cell that $p$ lies in are tested against each other to make
sure the fence post points lie on the opposite side of the implied
half plane cut defined by the normal of $(p-q)$ and will do so
as the fence grid grows.

Once the fence has been secured, extra points past the fence need to be
added to ensure that all nearby points that might affect connections to the
anchor point are included to the naive relative neighborhood graph calculation.
The $R _ {\text{max}}$ is calculated empirically during the course of
Algorithm \ref{alg:RNGpSPoIF} but is bounded above by a constant factor
(see Appendix \ref{TK}).

After fence securitization and all nearby relevant points have been collected,
a call to the naive relative neighborhood
graph, on the collected points relative to the anchor point, is called
in Algorithm \ref{alg:RNGPointNaive}.

The expected size of the secured fence is finite, in grid cell units,
and so the expected number of neighbors is almost surely bounded by
a constant.


\begin{algorithm}
  \caption{Single point Relative Neighborhood Graph, Secure Posts on Increasing Fence}\label{alg:RNGpSPoIF}
  \begin{algorithmic}[1]
    \Procedure{RNGPointSPoIF}{B, P, p}

      \State Initialize $\text{FencePostCluster}[]$ entries to $\text{False}$ %, \\
        %take $2 \times 2$ clusters of fence posts for each face

      \State $Q = \{\}, \ \ R _f = 0, \ \ R _ {\text{max}} = 0$
      %\State $Q = \{\}$
      \State $c = $ grid cell $p$ falls in

      \For{ $(R _ f \le N_s)$ \ {\bf and} \\  \hspace{3.0em} (\emph{False} $\in$ FencePostCluster) }
        \State $Q' = \text{GridPerimeterCellPoints}(B, p, R_f) / \{p\}$
        \For{ $q \in Q'$ }
          \State $Q = Q \cup \{q\}$
          \State $R' =$ max grid cell distance from $c$ \\
          \hspace{6.5em} $\text{Lune}(p,q)$ falls in
          %\State $R _ {\text{max}} = $ maximum of $R_{\text{max}}$ and grid cell distance that lune of $(p,q)$ intersects
          \State $R _ {\text{max}} = \max(R_{\text{max}}, R')$
        \EndFor

        \For{ $q \in  Q'$ }
          \State Define $\rho(u) = \text{sgn}((p - q) \cdot (u - q))$
          \State Define $\sigma(u) = \text{sgn}(u \cdot (u - q))$
          \For{ every cluster, $\gamma$ }
            \If{ $\gamma$ is out of bounds }
              \State $\text{FencePostCluster}[\gamma] = \text{\emph{True}}$
            \ElsIf{ every $v \in \gamma$, \\
                    \hspace{8.65em} $\rho(v) = -\rho(p)$ {\bf and } \\
                    \hspace{8.65em} $\sigma(v) = -\sigma(p)$ }
              \State $\text{FencePostCluster}[\gamma] = \text{\emph{True}}$
            \EndIf
          \EndFor
        \EndFor
        $R _ f = R _ f + 1 $
      \EndFor

      \For{ $R _ e = [R _ f+1 : R _ {\text{max}} ]$ }
        \State $Q'' = \text{GridPerimeterCellPoints}(B, p, R_e)$
        %\For{ $q \in \text{GridPerimeterCellPoints}(B, p, R_f) / \{p\}$ }
        \For{ $q \in Q''$ }
          \State $Q = Q \cup \{q\}$
        \EndFor
      \EndFor

      \State $\text{return } \text{RNGPointNaive}( Q, p )$

    \EndProcedure
  \end{algorithmic}
\end{algorithm}

%Each cluster in $\text{FencePostCluster}[]$ represents a small section of one of the bounding cube faces.
%One choice of clusters of points is a $3 \times 3$ grid of posts, evenly spaced on each of the bounding
%cube faces.
%The $3 \times 3$ posts are then grouped into neighbor $2 \times 2$ \emph{clusters}, for a total of $4$
%clusters per face ($6 \cdot 4 = 24$ in total, for all six faces).
%
%For a plane to partition a part of the bounding cube face, all posts within the cluster must be partitioned
%from the anchor point, which will be called \emph{securing} the cluster.
%Multiple clusters can be partitioned (\emph{secured}) by the same partioning plane, so long as individual
%clusters are fully partitioned.
%Clusters that fall out of bounds are marked as secured.
%Figure \ref{TK} shows an example highlighting the posts on a bounding cube face,
%its clusters, partitions and invalid partial partitions.

The auxiliary function, \emph{GridPerimeterCellPoints},
returns points within the grid, $B$,
centered at the grid cell that $p$ resides in and that lie an integral $R$ distance away.
The list of points within the enumeration the shell of cells at distance $R$ is returned, with cells
falling completely out of bounds skipped over.
A small note that the initial call of \emph{GridPerimeterCellPoints} needs to exclude anchor point,
$p$, as the starting grid cell will include the anchor point in question.

In the worst case, Algorithm \ref{alg:RNGpSPoIF} will iterate until it encompases all points within the grid
when the \emph{FencePostCluster}
has fallen completely out of bounds.
In such a case, Algorithm \ref{alg:RNGPointNaive} is run on all points and is equivalent to the degenerate
condition of running a naive $O(N^2)$ relative neighborhood graph computation on all neighbor points to an
anchor point.

For random point distribution in the unit cube, we expect much better runtime as the the expected
number of potential neighbors collected will be bounded.


\section{Motivation and Structure}

This section provides proofs that justify the algorithm's correctness and runtime.
Before going through specifics, an overview is provided, giving motivation
for the technical details that follow.

From the problem formulation, there are $N$ 3D points chosen uniformly at random in the unit cube.
As an initial pass,
the $N$ points are then binned into grid cells of size $N^{-1/3}$.

A single point is considered at a time
and the relative neighborhood graph connections are calcualted relative to this
point.
Each point considered in this way is called an \emph{anchor point}.
The runtime construction for the relative neighborhood graph relative to each
anchor point will be constant, providing an expected $O(N)$ algorithm once
each anchor point is processed.

To construct the relative neighborhood graph relative to the anchor point,
the algorithm makes use of a \emph{plane partition heuristic}.

For each anchor point, $p$, a nearby point, $q$, is used to create a half plane cut.
The half plane cut is defined from the plane that passes through $q$ with normal in
the direction of $(p-q)$, $H _ {q,p} = \{ u : (p - q) \cdot (u - q) > 0 \}$.
All points that fall outside of $H _ {q,p}$ can never have an edge to $p$ as the point $q$
would be in the lune excluding the connection.


\begin{figure}[!htb]
  \centering
  \includegraphics[width=0.9\linewidth]{rng_cutplane_heuristic.pdf}
  \label{fig:planePartitionHeuristic}
  \caption{ 2D representation of the plane partition heuristic. $p$ is the anchor point with $q$ a nearby
  point used to create the plane partition. The point $w$ is excluded by the plane partition heuristic as $q$ would
  be in the lune of $p$ and $w$. The point $u$ is not excluded by the plane partition heuristic, even if
  $q$ might fall within the lune created by $p$ and $u$ because $u$ falls below the plane partition. }
\end{figure}


Figure \ref{fig:planePartitionHeuristic} provides a 2D representation of the plane partition heuristic, showing
the anchor point $p$, a nearby  point $q$, a point that is excluded by the heuristic, $u$,
and a point, $w$, that still has a potential relative neighborhood graph connection to $p$.

If, for set of $q$ points, a compact convex hull can be created from the $H _ {q,p}$ half plane cuts,
then only points within the convex hull need be considered for relative neighborhood graph
connections to $p$.

Instead of working with the convex hull directly, a bounding region is used,
%The presented algorithm implicitly constructs a convex hull by considering a bounding region,
called a \emph{fence}, for which constructing a bounding convex hull is simpler.
The fence is a cube, centered at the grid cell center that the anchor point $p$ occupies.
The fence initially has side length of the cube cell spacing.
Each fence face is broken into square regions, called \emph{fence patches}, and each patch
is tested against the half plane cut.
In this paper, only nine fence posts per face are considered, with fence post patches
clustered into $2 \times 2$ posts, four clusters per fence face.

If the fence patch is completely outside of the half plane cut, the fence patch is said
to be \emph{secured}.
Figure \ref{TK} shows various patches that are in different states of securitization
from half plane cuts.

Points within the fence are used to construct the half plane cuts.
If the half plane cuts fail to secure the fence, the fence cube side is increased,
staying on grid cell boundaries.
The process is then repeated, attempting to secure the fence from the newly added points.
The fence is increased until it is secured, in the worst case increasing to the
degenerate condition that subsumes the whole unit cube.

Once all fence patches are secured, there is an implicit compact convex hull that is completely
within the fence cube and only points that fall within the fence need be considered
for relative neighborhood graph connections to the anchor point $p$.
Though only points inside of the fence cube have the possibility of a relative neighborhood
graph connection to $p$, points outside of the fence can still have an effect
by potentially restricting other connections to $p$.

The fence needs to be extended to add these additional points. As will be shown below,
the fence need only be extended by a constant proportion to the fence side length.
This adds an expected constant factor of points to consider, leaving the linear runtime
unaffected.

Once a fence has been secured and all relevant nearby points collected, a naive
relative neighborhood graph algorithm is then used to find connections from the anchor
point to its nerby  point collection.
The size of collected nearby points is expected linear, so the naive algorithm runs
in constant time.
%providing the expected linear time overall as each anchor point is considered in turn.

Figure \ref{TK} shows a graphical cartoon of this process as a summary.

The fence patch is reprsented by four points on its corners, called \emph{ fence posts}.
A fence post is secured if the patch, or \emph{cluster} of fence posts, falls
completely on the restricted side of a half plane cut.

As the algorithm operates by securing posts on an ever increasing fence until all
posts secured, the algorithm is named Securing Posts on Increasing Fence (SPoIF).

%The SPoIF algorithm is chosen for its conceptual simplicity 
%SPoIF has suffers from some tedious bookkeeping and details that need to be accounted for
%but provides a compromise
%of constant runtime efficiency and conceptual simplicity while maintaining the $O(N)$
%expected runtime.


% THIS GOES SOMEWHERE ELSE...CONCLUSION?

Proofs of correctness and algorithmic details are presented in section TK, TK.
Before proceeding with the technical aspects of the algorithm, a few points are discussed.

A secured fence provides a compact convex hull, the fence cube, and thus implies a smaller compact convex
hull bounded within it.
The converse is not true and a convex hull constructed from
an anchor point and its neighbors need not be secured by a fence in this way.
The convex hull could be considered directly, instead of the fence structure as proxy,
but this involves using convex hull planes defined by a point and normal with the need
to derive the vertex intersection points.
Since there are an expected constant number of half planes to consider, vertex enumeration conversion
would be constant but at the cost of an inflated runtime factor with a naive algorithm.
Using the convex hull directly is not considered further in this paper.

Below, nine fence posts per cube face (thirty six in total), will be used.
Some other number of fence posts could be used.
Another alternative is to use a shrinking area on each of the faces,
defined by the half plane intersections.

These are mentioned to note that many other methods could work and still retain
the expected linear time runtime.
Other than to mention them as possibilities, these ideas are not pursued further in this paper.

The nine fence posts per fence face implementation is now considered in detail.

\subsection{Overview}

In this section, we provide some technical proofs providing correctness for
the algorithm in section TK.
The general strategy will be to show that only nearby points within a compact
convex hull of an anchor point need be considered for relative neighborhood graph
connections, that nearby points within a convex factor distance of the convex hull
need to be considered for any influence and that the approximation by the coarse
fence and fence posts imply a contained compact convex hull.

Let $P = (p _ 0, p _ 1, \dots, p _ {n-1})$, where each $p _ i \in \mathbb{R}^3$ is
chosen uniformly at random in the unit cube.

% PROBABLY SHOULD BE TAKEN OUT
%...

A rough overview of the algorithm proceeds as follows:

\begin{itemize}
  \item { \bf Initialize} a three dimensional grid, where each cell size of $n^{-1/3}$ on a side
  \item { \bf Bin} points into grid cells, allowing for multiple points in a bin.
  \item For each point, $p$:
    \begin{itemize}
      \item {\bf Add} nearby points from adjacent grid cells
      \item {\bf Increase} the radius of a bounding cube, at grid cell increments,
        until the collected nearby points have precluded connections from
        all other points outside bounding cube to $p$
      \item {\bf Extend} the list of nearby points, once a bounding cube has been
        found, to include points that might preclude connections on the inner
        bounding cube
      \item {\bf Run} a naive relative neighborhood graph, \\
        $\text{RNG}(P,p)$, on all collected nearby points
        to determine connections to $p$
    \end{itemize}
\end{itemize}

The binning phase creates a Poisson process with an average of one point per bin.
A coarse grained plane partition heuristic is used to exclude further away points from consideration
and is the method by which the process knows to stop extending the bounding cube.

Once a bounding cube has been found, other points outside of the bounding cube, but still within
a, larger, local distance need to be added as they could influence connections on the interior bounding cube.
The list of all neighbors is given to a naive relative neighborhood graph algorithm and the process is then
repeated for each point.

Given points $p$ and $q$,  any point, $w$, that lies on the other side
of the plane with normal $\frac{(q-p)}{|q - p|}$ that passes through $q$, cannot connect to $p$, as $q$ would
be in any lune that $p$ and $w$ would create.
Figure \ref{TK} gives a graphical representation of points excluded by this plane heuristic
in two dimensions.

The entire region that $q$ excludes from $p$ is more complicated (see Figure \ref{TK})
so the plane partition heuristic is coarse but provides an
easy way to exclude large portions of points.

%The plane partition heuristic excludes points, but
To maintain a linear time algorithm, 
large collections of points that can never connect to our anchor point, $p$, need to be
efficiently excluded.
%we need to
%avoid processing points individually and
%to implicitely exclude large collection of points that can never connect to the our anchor point $p$.
To do this, consider a current grid cell collection of integral radius $R \in \mathbb{Z}$, centered
on the grid cell where anchor point $p$ resides.
If the tangent planes that passes through neighbor $q$, and in direction $(q-p)$ within the grid cell collection,
partition each of the six sides of our cube of radius $R$, we can be assured that any points outside
of the grid cell collection will never be able to connect to $p$.

A grid cell collection of a radius $R$ is called a \emph{fence} of radius $R \in \mathbb{Z}$.
When $R$ is understood from context, this will be called just a \emph{fence}.
If the sides of the \emph{fence} has been partitioned, with anchor point $p$ and neighbors $q$ within
the \emph{fence}, the \emph{fence} will be called \emph{secured}.

Points outside of the \emph{fence} will never be able to connect to anchor point $p$ but they still
might be in the lune of some points within the \emph{fence}.
For this reason, when collecting points for the naive relative neighborhood graph calculation, some
points outside of the fence will need to be added, as they can sabotage potential edges within
the \emph{fence}.

This will be done by considering a larger radius of the current \emph{fence} to add points
but the increase in $R$ will shown to be bounded by a constant factor, allowing overall linear time to be maintained.

Before discussing the algorithm in detail, preliminary structural results will be discussed.

\subsection{Foundation}

In this section we go through the proofs that provide the foundation of correctness
for the algorithm.


%Supporting mathematical foundation is provided in this section to motivate the
%algorithmic design.
%Our goal is to provide the mathematical infrastructure that considers connections
%to an \emph{anchor point} in a spatially local neighborhood around it.
%
%An overview of what's being achieved is:
%
%\begin{itemize}
%  \item Show the plane partition heuristic can be used to implicitely discard points
%  \item Show that the plane partition heuristic can be used to create convex hull region
%        of points that can connect
%  \item The convex hull need only be extended at most a constant proportion of the maximum
%        diameter to include the region that any point may reside in that can affect connections
%        to the anchor point
%  \item A bounding cube can be used as a proxy for convex hull, with the bounding cube
%        being partitioned implying a compact convex hull within it
%  \item The bounding cube, with high probability, need only be grown by
%        a constant factor, relative to grid cell size, ensuring constant runtime
%\end{itemize}
%
%Constant runtime for a finding relative neighborhood graph connections of a single
%point imply a linear runtime overal.
%
%From the anchor point to points nearby, a plane partition heuristic
%is used to partition space and to implicitely discard points from consideration.
%From the plane partition heuristic, a convex hull can be constructed
%whereby only points within need be considered for potential connections.
%Points outside the convex hull can't connect to the anchor point but
%still need to be considered as they still may be able to preclude connections.
%
%The convex hull will not be used directly.
%For conceptual simplicity and constant speed concerns,
%a bounding cube region, called a \emph{fence}, will be used.
%
%The bounding cube has points placed on each of its faces, called \emph{posts},
%grouped into clusters that will be used to indicate when the whole of the
%bounding cube is partitioned.
%The cube of grid cells is used to consider neighbor points and test
%plane partitions and is grown until the 
%
%A direct convex hull can be used but, for reasons discussed elsewhere, is
%not pursued further in this paper.

%\renewcommand\qedsymbol{$\blacksquare$}

%\newtheorem{theorem}{Theorem}

\newtheorem{theorem}{Theorem}
\newtheorem{lemma}[theorem]{Lemma}
\newdefinition{rmk}{Remark}
\newproof{pf}{Proof}
%\newproof{pot}{Proof of Theorem \ref{thmName}}


%\newtheorem{lemma}[theorem]{Lemma}

For an anchor point $p$, if $w$ is on the other side of the plane that passes through $q$,
with tangent in the direction of $(q-p)$,
then $q$ will always be in the lune created by $p$ and $w$ (see Figure \ref{fig:planePartitionHeuristic}):


\begin{lemma}[Plane Partition Heuristic]
  \label{lemma:plane}
  if $p,q,w \in \mathbb{R}^3$ with:
  %w \in H = \{ u \in \mathbb{R}^3 | (q-p) \cdot (u - q) \ge 0 \}
  $$
  \begin{array}{l}
  w \in H  _ {q,p} = \{ u \ :  \ (p-q) \cdot (u - q) > 0 \}
  \end{array}
  $$
  then
  $$
  \begin{array}{l}
    q \in \text{Lune}(p,w)  = \\
      \ \ \ \{ v \ : \ |v-p| \le |p-w| \ \& \ |v-w| \le |p-w| \}\\
  \end{array}
  $$
\end{lemma}

\begin{pf}
  TK
\end{pf}

The plane partition heuristic takes points near $p$ to be used in constructing a
convex hull around $p$.
Only nearby points that fall within the convex hull need to be tested for relative
neighborhood graph connections to $p$.

Lemma \ref{lemma:hull} states this more formally:

%Only points within the convex hull have the potential to connect to $p$:
%%need to be considered for a relative neighborhood graph connection to $p$:

\begin{lemma}[Convex Hull Heuristic]
  \label{lemma:hull}
  Let $p \in \mathbb{R}^3$, $Q = (q _ 0, q _ 1, \dots, q _ {m-1}) \subset P$ and $\mathbf{C}$ be the
  compact convex hull created from $H _ i = \{ u \in \mathbb{R}^3 : (p - q _ i) \cdot (u - q _ i) > 0 \}$,
  then $(w,p) \in \text{RNG}(P,p) \to w \in \mathbf{C}$
\end{lemma}

\begin{pf}
  By Lemma \ref{lemma:plane}, all $v$ outside of $\mathbf{C}$ will have some $q _ i \in Q$ in the $\text{Lune}(v,p)$, so cannot have an RNG edge to $p$.
\end{pf}

Lemma \ref{lemma:hull} only guarantees that RNG connections to $p$ fall within the hull, $\mathbf{C}$.
Not every point within $\mathbf{C}$ might connect to $p$ and points outside of $\mathbf{C}$ might still influence
connections to $p$.
For example, Figure \ref{TK} shows a potential connection between two points within the convex
hull but is thwarted by a point outside the convex hull that is within the lune created by the pair.

Points outside of the convex hall might influence connections to $p$ but we can bound how
far these points can be.
Nearby points influencing $p$ can only be a constant factor of the maximum distance of $p$ to the convex hull:

\begin{lemma}[Convex Hull Extension Heuristic]
  \label{lemma:extension}
  Let $p$ be the anchor point and $\mathbf{C}$ be the convex hull as detailed above.
  Let $R _ {\text{max}} = \max _ {u \in \mathbf{C}} |u-p|$ be the maximum distance of the convex hull boundary to $p$.
  All points that could affect $\text{RNG}(P,p)$ lie within a sphere of radius $R _ {\text{max}}$, centered at $p$.
\end{lemma}

\begin{pf}
  The maximum lune radius for a point, $u$, within the convex hull to $p$ is bounded above by $R _ {\text{max}}$.
  For a point, $w$, to influence a potential connection of $u$ to $p$, $w$ must lie within the lune created by
  $u$ and $p$.
  If $|w-p| > R _ {\text{max}}$, $w$ falls outside of any lune that $p$ might create with points within the convex
  hull.
\end{pf}

%%%%
%%%%
%%%%

% TAKE THIS OUT
% THIS DISCUSSION SHOULD GO IN CONCLUSION, NOT HERE
% WE KEEP MENTIONING IT, IT'S FINE, WE KNOW, JUST MENTION IT IN THE CONCLUSION

%The above already suggests a linear time algorithm by starting at a grid cell where $p$ resides
%and adding neighboring points from neighboring grid cells until a compact convex hull has been achieved
%that meets some maximum radius criteria.
%All points within a sphere centered at $p$ can then be added and a naive relative neighborhood graph algorithm
%can be run on the reduced neighbor set.
%It can be shown that for a constant number of neighbors, the maximum distance of $p$ to the convex hull is bounded,
%giving an expected linear time algorithm for random points in the unit cube.
%
%While there is an expected constant number of neighbors that make up the convex hull,
%a natural operation to find radius bounds would be to convert to
%a point representation via a vertex enumeration procedure, inflating the runtime constants were a
%naive vertex enumeration algorithm to be used.
%There may be ways of speeding up the vertex enumeration by considering the structure of the problem in question
%but this methodology is not pursued further in this paper.

%%%%
%%%%
%%%%

A bounding box cube is used as a proxy for the convex hull.
The bounding box cube, called a \emph{fence}, has each face partitioned into regular squares,
called \emph{fence face patches} (or just a \emph{fence patch}).
Plane partitions are used to partition the fence patch from the anchor point.
%has certain key points on each cube face that will be used
%to indicate when a plane has partitioned portions of the fence face from the anchor point, $p$.
Once all fence face patch portions have been partitioned from $p$, a process that will be called \emph{securing the fence},
this implies a compact convex hull contained within the fence.
%The fence cube is extended, as required by Lemma \ref{lemma:extension}, and all points within the extension
%are collected in consideration for the relative neighborhood graph connections to $p$.

%For the anchor point $p$, only $\text{RNG}(P,p)$ connections can happen within the secured fence but, as alluded
%to with \ref{lemma:extension}, other points outside of the secured fence might influence connections within.
%We will see that a constant factor of fence radius need be considered to add additional points to the naive
%relative neighborhood graph computation.

A secured fence necessarily implies a compact convex hull embedded within the fence but the other direction
is not true in general, as there can be a convex hull whose volume is completely contained within the
fence but whose partition planes do not secure the fence.
We will see that this has no effect on the expected linear runtime, as securing the fence for an individual
anchor point is expected to use a constant number of neighbors.

% probably a theorem...

% this needs to be split
% the patch secured in perpetuity condition is for the algorithm to not redo work after a patch has been secured
% the 

%\begin{theorem}[Cube Bounds on Interior Convex Hull]
\begin{theorem}
  \label{thm:cubeConvexHull}
  \ \\
  Let $\mathbf{B}$ an axis aligned cube centered at the origin, with side length $L _ {\mathbf{B}}$.
  Let each face be covered in $m$ non-overlapping axis aligned
  square patches, $\rho _ {d,j}$, where the four corners of patchs are
  $\rho _ {d,j,k}$ ($d \in \{ \pm x, \pm y, \pm z \}$, $j \in \{0 \dots (m-1)\}$, $k \in \{0,1,2,3\}$). \\
  %and with $N _ {\rho _ {d,j}}$ the normal of the patch away from the center. \\
  Let $\mathbf{H} = (H _ 0, H _ 1, \dots, H _ {s-1})$ be a list of half plane cuts with
  $Q = (q _ 0, q _ 1, \dots, q _ {s-1})$ be a respective point on the plane cut. \\
  If, for every patch $\rho _ {d,j}$ there exists a half plane $H _ i$, with normal $N _ {H _ i}$,
  such that every corner of the patch: \\ % $k \in \{0,1,2,3\}$,
  %$\rho _ {d,j,k} \cdot N _ {H _ i} < 0$ and
  $$
  (\rho _ {d,j,k} - q _ i) \cdot N _ {H _ i} < 0
  $$
  then $\cap _ i H _ i$ creates a compact convex hull that is strictly contained in $\mathbf{B}$.

  %for all axis aligned cubes, $\mathbf{B'}$ centered at zero of side length $L \ge L _ {\mathbf{B}}$
  %, for a given patch, all patch corner point are on the other
  %side of the half plane cut, $N _ {H _ i} \cdot (\rho _ {d,j,k} - q _ i) < 0$,
  %then
  %, $\cap _ i H _i$ ($\ne \emptyset$?) bounds a compact convex hull completely contained in $\mathbf{B'}$.
\end{theorem}

\begin{pf}
  The four points of $\rho _ {d,j}$, $(c\rho _ {d,j,0},\rho _ {d,j,1},\rho _ {d,j,2},\rho _ {d,j,3})$  create a frustum.
  A point lying on the patch is defined by $u = \sum _ k t _ k \rho _ {d,j,k}$ with $\sum _ k t _k = 1, t _k \ge 0$. \\
  Consider:\\
  $$
  \begin{array}{ll}
   & ( u - q _ i ) \cdot N _ {H _ i}  \\
    %= & ( (\sum _ k L t _ k \rho _ {d,j,k} ) - q _ i ) \cdot N _ {H _ i} \\
    %= & ( (\sum _ k L t _ k \rho _ {d,j,k} ) - (\sum _ k t _ k ) q _ i ) \cdot N _ {H _ i} \\
    %= & t _ k (\sum _ k (L-1) \rho _ {d,j,k} \cdot N _ { H _ i} ) + t _ k ( \rho _ {d,j,k} - q _ i ) \cdot N _ {H _ i} \\
    = & ( (\sum _ k t _ k \rho _ {d,j,k} ) - q _ i ) \cdot N _ {H _ i} \\
    = & ( (\sum _ k t _ k \rho _ {d,j,k} ) - (\sum _ k t _ k ) q _ i ) \cdot N _ {H _ i} \\
    = & \sum _ k [ t _ k ( \rho _ {d,j,k} - q _ i ) \cdot N _ {H _ i} ] \\
    < & 0
  \end{array}
  $$
  %Since $t _ k \ge 0$, $L \ge 1$ and each term is non-positive.
  Since $t _ k \ge 0$ with at least one non-zero $t _ k$ and $(\rho _ {d,j,k} - q_i) \cdot N _ {H _ i} < 0$ by assumption.
  Every point on the face of the cube $\mathbf{B}$ lies on the negative side of some half plane cut, bounding a compact
  convex hull within.
  %each point on the scaled patch, thus the whole patch, is partitioned by the plane defined by $H _ i$. \\
  %By assumption, all patches have been secured for $L = L _ {\mathbf{B}}$.
  %If there is a $w$ that was part of the convex hull but outside of the cube fence,
  %then $w$ must be on a patch for some fence size $L _ w$.
  %Since $L _ w > L$, this is contradiction, disallowing the existence of $w$.
\end{pf}

The SPoIF algorithm grows the fence if the fence patches haven't been secured.
In order to consider patches secured, an extra condition is needed to make sure the half plane will
keep the patch partitioned as the fence grows.


\begin{theorem}
  \label{thm:growSecure}
  \ \\
  Let $\mathbf{B}$, $\mathbf{H} = (H _ 0, H _ 1, \dots, H _ {m-1})$,
  $\mathbf{Q} = (q _ 0, q _ 1, \dots, q _ {s-1})$, $N _ {H _ i}$,
  $\rho _ {d,j}$, $\rho _ {d,j,k}$, 
  be as defined in Theorem \ref{thm:cubeConvexHull}. \\
  Let $\mathbf{B'}$ be an axis aligned cube centered at the origin with side length $L _ {\mathbf{B'}} \ge L _ {\mathbf{B}}$,
  $L = \frac{L _ {\mathbf{B'}}}{L _ {\mathbf{B}}} \ge 1$. \\
  If, for some patch $\rho _ {d,j}$, 
  \begin{itemize}
    \item[ ] $\rho _ {d,j,k} \cdot N _ {H _ i} < 0$ and 
    \item[ ] $(\rho _ {d,j,k} - q _ i) \cdot N _ {H _ i} < 0$
  \end{itemize}
  then the patch $L \cdot \rho _ {d,j}$ remains outside of the halfspace $H _ i$: $L \cdot \rho _ {d,j} \cap H _ i = \emptyset$.
\end{theorem}

\begin{pf}
  As in Theorem \ref{thm:cubeConvexHull}, the four points
  $\rho _ {d,j,k}$, ($k \in \{0,1,2,3\}$),
  with a point lying on the patch now defined by
  $v = \sum _ k L t _ k \rho _ {d,j,k}$ with $L \ge 1, \sum _ k t _k = 1, t _k \ge 0$. \\
  $$
  \begin{array}{ll}
   & ( v - q _ i ) \cdot N _ {H _ i}  \\
    = & ( (\sum _ k L t _ k \rho _ {d,j,k} ) - q _ i ) \cdot N _ {H _ i} \\
    = & ( (\sum _ k L t _ k \rho _ {d,j,k} ) - (\sum _ k t _ k ) q _ i ) \cdot N _ {H _ i} \\
    = & \sum _ k [ (L-1) t _ k \rho _ {d,j,k} \cdot N _ { H _ i}  + t _ k ( \rho _ {d,j,k} - q _ i ) \cdot N _ {H _ i} ] \\
    < & 0
  \end{array}
  $$
  %Since $t _ k \ge 0$, $L \ge 1$ and each term is non-positive.
  %Here, $L \ge 1$, $t _ k \ge 0$ with at least one non-zero $t _ k$.
  Both terms in the sum are negative by assumption of $\rho _ {d,j,k} \cdot N _ {H _ i} < 0$
  and $(\rho _ {d,j,k} - q_i) \cdot N _ {H _ i} < 0$, where the coefficients $(L-1)$ and $t _ k$
  are non-negative with at least one $t _ k$ positive.
\end{pf}

Theorem \ref{thm:cubeConvexHull} ensures that a compact convex hull in contained within the fence when the fence
patches have been partitioned.
Theorem \ref{thm:growSecure} ensures that once a patch has been secured, it remains secured
should the fence need to grow.
Theorem \ref{thm:growSecure} allows us to secure patches with candidate nearby points to an anchor point and not
reprocess the candidate points as they secure fence patches.
There may be optimizations that re-use candidate nearby points but these are not pursued further in this paper.



%Once a patch $\rho _ {d,j}$ is secured by a $H _ i$, it's secured in perpetuity for any larger bounding cube $\mathbf{B}$.
%Securing a patch is a stronger condition than having each of the patch points be partitioned for a given side length.
%A secure patch is partitioned no matter how big the fence cube gets, ensuring that as a cube fence grows, if a patch is
%secured at some point in the past, it will remain secure in the future.


%\begin{theorem}[RNGPointSPoIF($P$, $p$) runs in expected $O(1)$ time]
\begin{theorem}
  \label{thm:pspoifO1}
  The function RNGPointSPoIF($P$, $p$) runs in expected $O(1)$ time
\end{theorem}

\begin{pf}
  ...
\end{pf}

  Finding the relative neighborhood graph connections for $p$ is $O(1)$,
  providing us with the desired expected linear time algorithm:

%\begin{theorem}[RNGSPoIF($P$) runs in expected $O(N)$ time]
\begin{theorem}
  \label{thm:ppoifON}
  The function RNGSPoIF($P$) runs in expected $O(N)$ time
\end{theorem}

TODO:

- securing fence implies compact convex hull on interior

- a secured fence need only consider a larger fence, of factor $\sqrt{3}$, to include all points that could
  influence relative neighborhood edges to $p$

- $SPoIF:RNGp(P,p)$ runs in constant expected time

- $SPoIF:RNG(P)$ runs in expected linear time





\section{Conclusion}

An expected linear time algorithm to find the relative neighborhood graph for 3D points placed uniformly
at random in the unit cube.
To the author's knowledge, no well known expected linear time algorithm for finding the relative neighborhood
graph is known for the case of points distributed uniformly in the unit cube.

Randomness of point distribution eliminates many pathological cases that would need to be taken care of with
a fully deterministic algorithm.
The bounded maximum degree of the relative neighborhood graph in 3D, along with the randomness assumption, allow
a construction of an expected linear time algorithm by employing a coarse plane partition heuristic that allows
only local neighborhoods of points to be considered.
A finite neighborhood makes naive calculations constant time, allowing for an overall expected linear runtime algorithm.

The algorithm presented also answers a question posed by Katajainen and Nevalainen \cite{KN85} of whether there's a simpler
algorithm that wouldn't need to handle points close to the border differently.
The algorithm presented avoids needing to consider side or corner regions independently as the the plane partition
heuristic provides an orthogonal partition and fences near the edge of the region will be secured by falling out of
bounds.

Certain pathological simple geometries can cause problems for the algorithm as presented.
For example, a ``diamond" shaped distribution might have trouble securing the fence.
Extensions have the potential to overcome limitations while still retaining
overall linear time execution by,
%This can maybe be overcome by various methods,
for example, doing a pre-processing step
that marks grid cells as out of bounds if they fall outside of some initial convex hull calculation,
or doing random rotations and processing parallel copies.
These problems and ideas are not pursued in this paper other than to mention they exist and there might be
a way to mitigate them.






\section{Bibliography}

%Two bibliographic style files (\verb+*.bst+) are provided ---
%{model1-num-names.bst} and {model2-names.bst} --- the first one can be
%used for the numbered scheme. This can also be used for the numbered
%with new options of {natbib.sty}. The second one is for the author year
%scheme. When  you use model2-names.bst, the citation commands will be
%like \verb+\citep+,  \verb+\citet+, \verb+\citealt+ etc. However when
%you use model1-num-names.bst, you may use only \verb+\cite+ command.
%
%\verb+thebibliography+ environment.  Each reference is a\linebreak
%\verb+\bibitem+ and each \verb+\bibitem+ is identified by a label,
%by which it can be cited in the text:
%
%\noindent In connection with cross-referencing and
%possible future hyperlinking it is not a good idea to collect
%more that one literature item in one \verb+\bibitem+.  The
%so-called Harvard or author-year style of referencing is enabled
%by the \LaTeX{} package {natbib}. With this package the
%literature can be cited as follows:
%
%\begin{enumerate}[\textbullet]
%\item Parenthetical: \verb+\citep{WB96}+ produces (Wettig \& Brown, 1996).
%\item Textual: \verb+\citet{ESG96}+ produces Elson et al. (1996).
%\item An affix and part of a reference:\break
%\verb+\citep[e.g.][Ch. 2]{Gea97}+ produces (e.g. Governato et
%al., 1997, Ch. 2).
%\end{enumerate}
%
%In the numbered scheme of citation, \verb+\cite{<label>}+ is used,
%since \verb+\citep+ or \verb+\citet+ has no relevance in the numbered
%scheme.  {natbib} package is loaded by {cas-dc} with
%\verb+numbers+ as default option.  You can change this to author-year
%or harvard scheme by adding option \verb+authoryear+ in the class
%loading command.  If you want to use more options of the {natbib}
%package, you can do so with the \verb+\biboptions+ command.  For
%details of various options of the {natbib} package, please take a
%look at the {natbib} documentation, which is part of any standard
%\LaTeX{} installation.

\appendix
\section{Appendix.0}

For completeness, this section discusses the elementary, but tedious, calculations
for the dart shape and volume.
As a reminder, our goal is to show that increasing the grid fence by a fixed size
will almost surely encapsulate a compact convex hull made from an anchor point and its neighbors.
The convex hull guarantees that the relative neighborhood graph edges of the anchor point to
its neighbors need only consider candidate points within the convex hull (with some caveats
about extending the grid fence to include saboteurs, detailed in section \ref{TK}).

The logic chain is roughly:

\begin{itemize}
  \item For an anchor point, $p$, and a grid fence side, there will be a region that a neighbor
        point can fall in that will partition the fence side from the point, a process we call
        \emph{securing the fence}, where a secured fence necessarily implies a compact convex hull
        made from the points on the interior
  \item The volume that secures a fence side will turn out to be four sphere slices joined together,
        which we will call a \emph{dart}, where a sphere slice is the intersection of two perpendicular
        half planes and a sphere
  \item A lower bound estimate on the dart volume is a constant proportion of overall grid fence volume
  \item A non-trivial constant dart volume makes the question of how many neighbors to be expected
        to secure the overal six sided fence into a coupon collector-like problem, giving us an
        expected finite grid fence size
  \item Finally, approximatations on the four slices bundled together
        allow us to give coarse values for grid fence sizes
\end{itemize}

It should be stressed that the discussion and estimates are meant primarily to show existence of
a finite grid fence size.
We focus on securing a fence side by waiting for some single neighbor point to secure it.
%, which is equivalent to considering a four post fence side.
The Shrinking Posts on Increasing Fence (SPoIF) algorithm in this paper has a smaller expected grid fence
size before securitization
as points can secure parts of the fence through clusters of posts, increasing the admissible
region on each face side that can be secured.

The calculations done in this appendix are meant to show that we only need to consider a finite grid fence
size by providing correct, if horrible, bounds.
In this paper, we are satisfied with providing empherical results for grid sizes that SPoIF uses without
going into more rigorous analysis.



First consider a fixed point $q$ and a fixed but moveable point $\eta$.

We think of $\eta$ as a point on a plane with normal $\eta / |\eta|$
that intersects $q$ in the slice of $z=0$ to $z=1$.

As we vary $\eta$ we want to know what the shape is traced out.

The claim is that this is the surface of a sphere whose center is at $q/2$
of radius $q/2$.

To see this, consider $\eta = q/2 + (|q| / 2) v$, with $|v| = 1$.

Calculate dot product of the line from $q$ to $\eta$ with the line from the origin to $\eta$:

$$
\begin{array}{ll}
  & \eta = q/2 + (|q| / 2) v \\
  & (q/2 + (|q|/2) v) \cdot (q/2 + (|q|/2) v - q ) \\
  = & (q/2 + (|q|/2) v) \cdot ( -q/2 + (|q|/2) v ) \\
  = & -(1/4)|q|^2 + (|q|/2) ^2 |v|^2 \\
  = & 0
\end{array}
$$

Proving the claim.

We want to find the the minimum volume of the dart
for some point inside of a fence of radius $s$ (side length $2s$)
nested inside of a larger fence of radius $s(2k+1)$ (side length $2s(2k+1)$).

First consider a sphere of radius $R$, centered at \\
$(-D,-D,-D)$ intersected with two
half planes, $H ^ + _ {x=0} = \{ u : u _ x \ge 0 \}$, $H ^ + _ {y=0} = \{ u : u _ y \ge 0 \}$.
For suitable $R$ and $D$, we want to estimate the volume.

If we take two tetrahedrons completely enclosed by the slice, the maximum diagonal of
the slice is $(R-\sqrt{2}D)$, with maximum straight side length (in the $x$-axis line and $y$-axis line) to
also be $(R-\sqrt{2}D)$ and slice half height to also be $(R-\sqrt{2}D)$.

The volume of the tetrahedron is then:

$$
(1/6) (R - \sqrt{2} D)^2 \sqrt{R^2 - 2 D^2}
$$

The volume of the two tetrahedrons stacked on top of each other:

$$
(1/3) (R - \sqrt{2} D)^2 \sqrt{R^2 - 2 D^2}
$$

More succinctly:

$$
\begin{array}{l}
  V _ {\text{slice}} > \frac{1}{3} (R - \sqrt{2} D)^2 \sqrt{R^2 - 2 D ^ 2}
\end{array}
$$

We will get concrete estimates below but as a quick confirmation, we can check to make sure that
this slice is a constant proportion of the fence it sits in.

We know that $R \propto (k s)$
and $D \propto (k s)$, for some $k \in \mathbb{Z}$ and grid cell size $s \in \mathbb{R}$.
Call $R = \alpha (k s)$ and $D = \beta (k s)$.

The dart will be four slices bundled together sitting inside of a $\gamma^3 (k s)^3$ volume.
Take $\alpha, \beta \in \mathbb{R}$ to be the smallest of the four slices:

$$
\begin{array}{ll}
  & \frac{ (4/3) (\alpha (k s) - \sqrt{2} \beta (k s))^2 \sqrt{ \alpha^2 (k s)^2 - 2 \beta^2 (k s)^2 } }{ \gamma^3 (k s)^3 } \\[10pt]
  = & \frac{ (4/3) (k s)^3 (\alpha - \sqrt{2} \beta )^2  \sqrt{ \alpha^2 - 2 \beta^2 } }{ \gamma^3 (k s)^3 } \\[10pt]
  = & \frac{ (4/3) (\alpha - \sqrt{2} \beta )^2  \sqrt{ \alpha^2 - 2 \beta^2 } }{ \gamma^3 } \\[10pt]
\end{array}
$$

as desired.

When choosing random points inside of the larger fence cube, we can think of the dart as the probability
a single point secures one side of the fence.
It is equivalent to talk about securing one side of the fence face with securing four posts on each of corners of
the fence face.
As a reminder, the Shrinking Posts on Increasing Fence (SPoIF) uses nine posts per face, equally distributed
along the face square.
Since there are more opportunities to secure clusters of posts on a fence face,
in addition to still retaining the features of the dart region, SPoIF will be able to secure the fence cube
more quickly.
We only consider the four post fence here to help simplify the calculation even though it inflates the expected
number of points we need to see to secure the grid fence.

Take the minimum dart value volume of the six directions and call it $V _ {\text{dart}}$.
The probability that a point lands in one of the darts is bounded below by $V _ {\text{dart}} / V _ {\text{fence}}$.
Of the points that land in one of the darts, the time we need to take to secure each face of the fence is now
a coupon collector-like problem with six coupons.

We can now plug in some real numbers to get coarse bounds on dart volumes, probabilities of landing in the critical region
and expected grid sizes.

As a reminder, we are considering a center grid cell, of side length $s$ where the anchor point, $p$,
resides, all of which is inside a larger grid fence of side length $2s(k+1)$, for some fixed $k \ge 0, k \in \mathbb{Z}$
as yet to be determined.

The smallest size a slice can be is if the anchor point is on the corner of the interior grid cell and to choose the grid fence corner
nearest to it.
For example, $p = (s/2,s/2,s/2)$, $q = ( s(k+1/2), s(k + 1/2), s(k+1/2) )$.

In such a case, the center point of the sphere is:

$$
(p + q)/2 = ( s(k+1)/2, s(k+1)/2, s(k+1)/2 )
$$

with radius:
$$
\begin{array}{ll}
  R & = |(p+q)/2 - p| \\
   & = |(sk/2,sk/2,sk/2)| \\
   & = (\sqrt{3}/2) sk
\end{array}
$$

Here:

$$
\begin{array}{ll}
  V _ {\text{fence}} & = (s(2k+1))^3 \\
  D & = s(k+1)/2 - s/2 \\
  & = sk / 2 \\
\end{array}
$$

This yields:

$$
\begin{array}{ll}
  \frac{ V_{\text{dart}} }{ V _ {\text{fence}}} > &  \frac{ \frac{4}{3} ( \frac{\sqrt{3}}{2} sk - \frac{\sqrt{2}}{2} sk)^2 ( \frac{3}{4}s^2 k^2 - \frac{1}{2} s^2 k^2 )^{\frac{1}{2}}  }{  s^3 (2k+1)^3  } \\[10pt]
  = & \frac{ s^3 k^3 ( \sqrt{3} - \sqrt{2} )^2 }{ 3 \sqrt{2} s^3 (2k+1)^3 } \\[10pt]
  = & \frac{ ( \sqrt{3} - \sqrt{2} )^2 }{ 3 \sqrt{2} (2+\frac{1}{k})^3 } \\[10pt]
  > & \frac{ ( \sqrt{3} - \sqrt{2} )^2 }{ 81 \sqrt{2}  } \ \ \ \ \ \ \ \ (k \ge 1) \\[10pt]
  > & 1 / 1198
\end{array}
$$

For points that land within one of the six dart region, we would expect this to be be a six coupon collector
problem with, at worst, equal probability.
The coupon collector problem grows roughly as $n \log(n)$  which, in our case of six sides of the fence to secure,
gives us an expected $11$ or more points before we secure the fence.

Each point has roughly a $6 / 1198$ chance of landing in one of the six darts, so we would expect
roughly $11 / (6 / 1198) = 2196.3\dots$ number of points within the grid fence to appear before
securing the fence.

Each grid cell has roughly one point, so we would expect the number grid cells on a side to
be

$$
\begin{array}{ll}
  & (11 \cdot 1198 / 6)^{\frac{1}{3}} \\
  & = 12.9\ldots \\
  & \sim 13 \\
\end{array}
$$

That is, we would expect $k \sim 6$, iterating roughly six to seven steps of increasing the grid fence size
until securing the fence in this way.


\section{Appendix.1}
...
Appendix sections are coded under \verb+\appendix+.

\printcredits

%% Loading bibliography style file
%\bibliographystyle{model1-num-names}
\bibliographystyle{cas-model2-names}

% Loading bibliography database
\bibliography{cas-refs}


%\vskip3pt

\bio{}
...
\endbio

\end{document}

